% !PRINTED 141
\weberpage{141}

\begin{center}
\textbf{\S\,43. \gesperrt{Die Jacobischen Funktionen} $\Theta(v)$, $H(v)$.}
\end{center}

Wenn man die $\vartheta$-Funktionen der Variablen $u$ als Funktionen
von $v$ betrachtet, so erhält man die \gesperrt{Jacobischen} $\Theta$-Funktionen.
Wir setzen
\weq{1}{\begin{aligned}
\vartheta_{01}\left(\frac{v}{2K}, \omega\right) &= \Theta(v),\\
\vartheta_{11}\left(\frac{v}{2K}, \omega\right) &= H(v).
\end{aligned}}
Es sind dann $\Theta(v)$, $H(v)$ als Funktionen $t(v, 2K, 2K')$ zu bezeichnen, und nach \S\,25 ist
\weq{2}{\begin{aligned}
\Theta(v) &= 1 - 2q\cos\frac{\pi v}{K} + 2q^4\cos\frac{2\pi v}{K} - 2q^9\cos\frac{3\pi v}{K} + \cdots\\
H(v) &= 2q^{\frac{1}{4}}\sin\frac{\pi v}{2K} - 2q^{\frac{9}{4}}\sin\frac{3\pi v}{2K} + 2q^{\frac{25}{4}}\sin\frac{5\pi v}{2K} - \cdots\\
q &= e^{-\frac{\pi K'}{K}}.
\end{aligned}}
Nach \S\,21 (8) ist
\weq{3}{\begin{aligned}
H(v) &= -ie^{-\frac{\pi K'}{4K}+\frac{\pi iv}{2K}}\Theta(v+iK'),\\
\Theta(v) &= -ie^{-\frac{\pi K'}{4K}+\frac{\pi iv}{2K}}H(v+iK').
\end{aligned}}
Ferner:
\weq{4}{\begin{aligned}
\vartheta_{00}\left(\frac{v}{2K}, \omega\right) &= \Theta(v+K),\\
\vartheta_{10}\left(\frac{v}{2K}, \omega\right) &= H(v+K),
\end{aligned}}
wonach, mittels der Formeln (3), (10), (15) des vorigen Paragraphen und (5), \S\,23:
\[
\sqrt{\varkappa} = \frac{H(K)}{\Theta(K)}, \qquad \sqrt{\varkappa'} = \frac{\Theta(0)}{\Theta(K)},
\]
\weq{5}{\begin{aligned}
&\Theta(K) = \sqrt{\frac{2K}{\pi}}, \quad \Theta(0) = \sqrt{\frac{2\varkappa'K}{\pi}}, \quad H(K) = \sqrt{\frac{2\varkappa K}{\pi}},\\
&H'(0) = \frac{\Theta(0)H(K)}{\Theta(K)} = \sqrt{\frac{2\varkappa\varkappa'K}{\pi}};
\end{aligned}}
endlich die Darstellung der elliptischen Funktionen:
%CONT>>
